\documentclass{article}
\usepackage{anyfontsize}
\usepackage[UTF8]{ctex} % 如果需要支持中文
\usepackage{fontspec} 
\setCJKmainfont{SourceHanSansCN-Light}[
    BoldFont=SourceHanSansCN-Medium
]

\usepackage[a4paper, left=0.1in, right=0.1in, top=0.05in, bottom=0.05in]{geometry} % 设置四边页边距为0.1英寸{geometry} % 设置页边距为0.5英寸
\usepackage{enumitem} % 用于自定义列表
\usepackage{amsmath}  % 数学公式支持
\usepackage{tikz}
\usetikzlibrary{arrows.meta, positioning}
\setlength{\parindent}{0pt} % 不缩进
\usepackage{etoolbox} % 用于列表操作
%调整类代码
\usepackage{comment}
\usepackage{tocloft} % 用于定制目录 样式
\usepackage{hyperref} %不需要目录盒子注释这
\usepackage{ifthen}
\usepackage{tikz}
\usetikzlibrary{arrows.meta}
\usepackage{titletoc}
\usepackage{graphicx}
\usepackage{float}

\usepackage{amssymb}  % 提供数学符号，包括\therefore
%目录设定
%快捷键 CMD + / 注释
%  \renewcommand{\cftdot}{} % 隐藏目录中的页码
%  \cftpagenumbersoff{section}
%  \cftpagenumbersoff{subsection}
%  \makeatletter
%  \renewcommand{\@pnumwidth}{0em}  % 设置页码宽度为0
%  \renewcommand{\@tocrmarg}{0em}   % 设置右边距为0
%  \makeatother
 


\usepackage{environ}

\newif\ifshowcontent  % 定义一个条件判断变量
\showcontenttrue    % 设置为 false，隐藏所有 question 环境的内容

\newcounter{question} % 题目计数器
\setcounter{question}{0}
\NewEnviron{question}{%
    \ifshowcontent
        \par\stepcounter{question}\textbf{\thequestion.} \BODY\par % 如果显示内容，则输出
    \fi
}

\NewEnviron{choices}{%
    \ifshowcontent
        \begin{enumerate}[label=\Alph*., leftmargin=*]
            \BODY
        \end{enumerate}\par
    \fi
}

\NewEnviron{correctanswer}{%
    \ifshowcontent
        \textbf{答案:}\BODY\par
    \fi
}



\newenvironment{explanation}
{\textbf{解析:}\par}
{\par}



% 新增考察方面命令
\newcommand{\checklist}[1]{
    \textbf{考察:} #1 \par % 在题目下方显示考察方面
    \addcontentsline{toc}{subsection}{题号\thequestion 考察: #1} % 将考察内容添加到目录
    \vspace{2em} % 添加1em的垂直空间
}

\graphicspath{{images/}} %统一


\begin{document}
 %\fontsize{22}{31}\selectfont
 \tableofcontents % 添加目录
\newpage
%\pagestyle{empty}




\section{考点1:信用风险建模评估}
\setcounter{question}{0}

\begin{question}
    Which of the following statements regarding a banking credit analyst's skills is most likely correct?
\end{question}

\begin{choices}
    \item High earnings quality suggests that the bank is profitable
    \item Peer analysis is facilitated by the standardized nature of financial performance measures
    \item Although qualitative analytical skills are required, quantitative analytical skills are more important
    \item In analyzing an unfamiliar banking sector, an analyst should start by performing detailed reviews of the major banks
\end{choices}

\begin{correctanswer}
    B
\end{correctanswer}

\begin{explanation}
    \textbf{A选项错误。}
    \begin{itemize}
        \item 高收益质量并不一定意味着银行是盈利的。收益质量指的是报告收益的可靠性和一致性，而不是盈利能力。
    \end{itemize}
    
    \textbf{B选项正确。}
    \begin{itemize}
        \item 同业分析是指将目标银行与类似银行和金融机构进行比较（财务和信用度）。由于财务业绩指标具有标准化性质，这使得同业分析变得更加容易。
    \end{itemize}
    
    \textbf{C选项错误。}
    \begin{itemize}
        \item 定量和定性分析技能同样重要，它们服务于不同（但相关）的目的。定性技能有助于确定实体偿还债务的意愿，而定量技能有助于确定实体偿还债务的能力。
    \end{itemize}
    
    \textbf{D选项错误。}
    \begin{itemize}
        \item 在分析一个不熟悉的银行业时，分析师应该从对监管的总体结构、特征和性质的初步研究开始。之后，才能对最大（其次是较小的）银行进行详细审查。
    \end{itemize}
\end{explanation}

\checklist{银行信用分析的基本概念和方法}



\begin{question}
    An exposure has a stressed default probability of 4.0\% and loss given default of 50.0\%. The portfolio with a value of \$1 million. The standard deviation of the LGD is 25.0\%. What is the stressed expected loss?
\end{question}

\begin{choices}
    \item 0.02 million
    \item 0.0125 million
    \item 0.04 million
    \item 0.05 million
\end{choices}

\begin{correctanswer}
    A
\end{correctanswer}

\begin{explanation}
    \textbf{步骤1：确定参数}
    \begin{itemize}
        \item 组合价值（EAD）：\$1{,}000{,}000
        \item 压力违约概率（PD）：4.0\% = 0.04
        \item 违约损失率（LGD）：50.0\% = 0.50
        \item LGD的标准差：25.0\%（计算预期损失时不需要，因为预期损失只使用期望值）
    \end{itemize}

    \textbf{步骤2：应用预期损失公式}
    \begin{itemize}
        \item 预期损失公式：
        \[
        EL = EAD \times PD \times LGD
        \]
    \end{itemize}

    \textbf{步骤3：计算压力预期损失}
    \begin{itemize}
        \item 代入数值：
        \[
        EL = 1{,}000{,}000 \times 0.04 \times 0.50 = 20{,}000 = 0.02 \text{ million}
        \]
    \end{itemize}

 
    
\end{explanation}

\checklist{压力下的预期损失计算}


\begin{question}
    Which of the following statements regarding the role of a corporate credit analyst is most likely correct?
\end{question}

\begin{choices}
    \item Earnings analysis is by far the most important analyst task
    \item The larger the size of the firm, the lower the cost of analysis
    \item Analysts are generally required to cover multiple industry areas given the huge diversity among corporations
    \item The smaller the firm, the lower the cost of analysis
\end{choices}

\begin{correctanswer}
    B
\end{correctanswer}

\begin{explanation}
    \textbf{A选项错误。}
    \begin{itemize}
        \item 现金流分析，而不是收益分析，是评估企业信用风险的关键。收益分析只是信用分析的一个组成部分。
    \end{itemize}
    
    \textbf{B选项正确。}
    \begin{itemize}
        \item 对于大型上市公司，可能有大量公开信息可供参考，只需要进行二次研究，从而降低分析成本。而对于较小的私营公司，可用的信息较少，因此需要更多的尽职调查和初步研究，从而增加分析成本。
    \end{itemize}
    
    \textbf{C选项错误。}
    \begin{itemize}
        \item 尽管基本的分析原则相同，但被分析公司在业务部门、产品、规模和地理位置方面存在巨大差异。因此，公司信用分析师必须拥有特定的行业知识才能有效工作。分析师很可能只专注于一个或两个行业领域。
    \end{itemize}
    
    \textbf{D选项错误。}
    \begin{itemize}
        \item 较小的公司通常信息较少，需要更多的尽职调查和初步研究，因此分析成本更高，而不是更低。
    \end{itemize}
\end{explanation}

\checklist{公司规模与信用分析成本的关系}


\begin{question}
    Which of the following is true about predatory lending and predatory borrowing?
\end{question}

\begin{choices}
    \item Both underprovide credit
    \item Both overprovide credit
    \item Predatory lending underprovides credit and predatory borrowing overprovides credit
    \item Predatory lending overprovides credit and predatory borrowing underprovides credit
\end{choices}

\begin{correctanswer}
    B
\end{correctanswer}

\begin{explanation}
    \textbf{A选项错误。}
    \begin{itemize}
        \item 掠夺性借出和掠夺性借入都不是授信不足，而是授信过度。
    \end{itemize}
    
    \textbf{B选项正确。}
    \begin{itemize}
        \item 掠夺性借出（Predatory Lending）指的是借款人不一定能申请到这笔贷款，或者不是必须要进行融资，但是由于金融机构强烈建议其融资，使得借款人借入与自身条件不匹配的贷款。
        \item 掠夺性借入（Predatory Borrowing）指的是借款人通过伪造信息借到了本身无法获得的贷款，比如与房产经纪商、抵押贷款经纪商、估价师和律师等进行勾结。
        \item 这两种情况都属于过度授信，都会导致信用风险增加。
    \end{itemize}
    
    \textbf{C选项错误。}
    \begin{itemize}
        \item 掠夺性借出是过度授信，而不是授信不足。掠夺性借入也是过度授信，而不是授信过度。
    \end{itemize}
    
    \textbf{D选项错误。}
    \begin{itemize}
        \item 掠夺性借出是过度授信，而不是授信过度。掠夺性借入也是过度授信，而不是授信不足。
    \end{itemize}
\end{explanation}

\checklist{掠夺性借贷的定义}


\begin{question}
    Which of the following model calculates the change in portfolio value due to rating migration of the underlying instruments?
\end{question}

\begin{choices}
    \item Credit Risk+
    \item Credit Metrics
    \item KMV
    \item Both a and c above are true
\end{choices}

\begin{correctanswer}
    B
\end{correctanswer}

\begin{explanation}
    \textbf{A选项错误。}
    \begin{itemize}
        \item Credit Risk+模型主要关注违约风险，而不是评级迁移。它使用泊松分布来模拟违约事件，不直接考虑评级变化对投资组合价值的影响。
    \end{itemize}
    
    \textbf{B选项正确。}
    \begin{itemize}
        \item CreditMetrics模型专门用于计算由于基础债券的信用迁移（如信用利差的变化）而导致的投资组合价值的变化。
        \item 该模型通过模拟评级迁移矩阵来评估信用风险，并计算相应的投资组合价值变化。
    \end{itemize}
    
    \textbf{C选项错误。}
    \begin{itemize}
        \item KMV模型主要关注违约概率的估计，使用期权定价理论来评估违约风险，而不是直接计算评级迁移对投资组合价值的影响。
    \end{itemize}
    
    \textbf{D选项错误。}
    \begin{itemize}
        \item 由于A和C选项都不正确，所以这个选项也是错误的。
    \end{itemize}
\end{explanation}

\checklist{信用迁移对投资组合价值的影响}

\begin{question}
    Lisa Chen is a new loan officer at Pacific Bank. Her manager requires her to make five commercial loans this month. Chen hears about a local entrepreneur with a good reputation. She meets with the entrepreneur, David Wang, and is impressed by his business acumen. They discuss a loan for Wang's new project, and Chen agrees it's a promising opportunity. She convinces the loan committee that Pacific Bank is fortunate to work with someone of Wang's reputation. This is an example of:
\end{question}

\begin{choices}
    \item Historical analysis technique
    \item Qualitative analysis technique
    \item Quantitative analysis technique
    \item Extrapolation analysis technique
\end{choices}

\begin{correctanswer}
    B
\end{correctanswer}

\begin{explanation}
    \textbf{A选项错误。}
    \begin{itemize}
        \item 历史分析技术主要关注过去的财务数据和表现，而本题中并没有涉及对历史数据的分析。
    \end{itemize}
    
    \textbf{B选项正确。}
    \begin{itemize}
        \item 名义贷款（Name lending）是一种定性分析技术，有时用来代替财务分析。
        \item 这种技术主要用于评估借款人偿还债务的意愿，而不是通过具体的财务数据来评估。
        \item 在本题中，Chen主要基于Wang的声誉和商业头脑来做决策，而不是基于具体的财务分析。
    \end{itemize}
    
    \textbf{C选项错误。}
    \begin{itemize}
        \item 定量分析技术主要关注具体的财务数据和指标，而本题中并没有涉及对具体财务数据的分析。
    \end{itemize}
    
    \textbf{D选项错误。}
    \begin{itemize}
        \item 外推分析技术主要基于历史数据进行预测，而本题中并没有涉及对历史数据的分析和预测。
    \end{itemize}
\end{explanation}

\checklist{名义贷款的定义}


\begin{question}
    A bank has an outstanding trade with one of its counterparties with an exposure of \$500,000 and a recovery rate of 70\%. The bank estimated that there is a 2\% probability that the counterparty will default on its obligations. What is the bank's expected loss?
\end{question}

\begin{choices}
    \item \$3,000
    \item \$7,000
    \item \$10,000
    \item \$150,000
\end{choices}

\begin{correctanswer}
    A
\end{correctanswer}

\begin{explanation}
    \textbf{步骤1：确定参数}
    \begin{itemize}
        \item 风险敞口（EAD）：\$500{,}000
        \item 回收率（RR）：70\% = 0.70
        \item 违约概率（PD）：2\% = 0.02
    \end{itemize}

    \textbf{步骤2：计算违约损失率（LGD）}
    \begin{itemize}
        \item LGD = 1 - RR = 1 - 0.70 = 0.30 = 30\%
        \item 回收金额 = \$500{,}000 × 70\% = \$350{,}000
        \item LGD = \$500{,}000 - \$350{,}000 = \$150{,}000（绝对金额）
    \end{itemize}

    \textbf{步骤3：计算预期损失}
    \begin{itemize}
        \item 预期损失公式：
        \[
        EL = EAD \times PD \times LGD
        \]
        \item 代入数值：
        \[
        EL = 500{,}000 \times 0.02 \times 0.30 = 500{,}000 \times 0.006 = \$3{,}000
        \]
    \end{itemize}
    
\end{explanation}

\checklist{预期损失的计算}

\begin{question}
    Michael Zhang, a credit analyst at City Bank, is reviewing a loan application from a local car dealership. The dealership has been owned by David Lee for over 20 years and sells American cars. Due to its suburban location, the dealership mainly sells SUVs and trucks. However, sales have declined, and fuel prices have risen. Lee is considering adding hybrid cars to his inventory. Lee has borrowed from City Bank before but currently has no outstanding loans. Which of the following statements is not one of the four components of credit analysis Zhang should be evaluating for this potential loan?
\end{question}

\begin{choices}
    \item The business environment, competition, and economic climate in the region
    \item Lee's character and past payment history with the bank
    \item The car dealership's balance sheets and income statements for the last few years as well as Lee's personal financial situation
    \item The financial health of Lee's friends and family who could be called upon to guarantee the loan
\end{choices}

\begin{correctanswer}
    D
\end{correctanswer}

\begin{explanation}
    \textbf{A选项错误。}
    \begin{itemize}
        \item 这是信用风险评估的第二个要素：外部环境（External Conditions）。包括国家风险和经营条件等，如当地的经济环境和竞争状况。
    \end{itemize}
    
    \textbf{B选项错误。}
    \begin{itemize}
        \item 这是信用风险评估的第一个要素：债务人的偿付能力和还款意愿（Obligor's Capacity and Willingness to Repay）。包括债务人的信用质量和过去的还款记录。
    \end{itemize}
    
    \textbf{C选项错误。}
    \begin{itemize}
        \item 这是信用风险评估的第一个要素：债务人的偿付能力和还款意愿。包括企业的财务状况和债务人的个人财务状况。
    \end{itemize}
    
    \textbf{D选项正确。}
    \begin{itemize}
        \item 信用风险评估的四个主要要素是：
        \begin{enumerate}
            \item 债务人的偿付能力和还款意愿
            \item 外部环境
            \item 特定信用工具的相关特征
            \item 信用风险缓释工具（如抵押品和贷款担保）
        \end{enumerate}
        \item 虽然风险缓释工具是信用分析的一部分，但对于一个经营20年的本地汽车经销商来说，不太可能寻求朋友或家人的贷款担保。即使需要担保人，分析师也应该评估具体的担保人，而不是简单地评估"朋友和家人"。
    \end{itemize}
\end{explanation}

\checklist{信用风险评估的四个主要要素}

\begin{question}
    The KMV model measures the normalized "distance to default". How is this defined?
\end{question}

\begin{choices}
    \item (Expected Assets - Weighted Debt) / (Volatility of assets)
    \item Equity / (Volatility of equity)
    \item Probability of stock price falling below a threshold
    \item Leverage x Stock Price Volatility
\end{choices}

\begin{correctanswer}
    A
\end{correctanswer}

\begin{explanation}
    \textbf{A选项正确。}
    \begin{itemize}
        \item 违约距离（Distance to Default）的定义是：（预期资产 - 加权债务）/（资产波动率）。
        \item 这是一个标准化的违约衡量指标，因此可以用于比较不同公司。
        \item 它本身并不直接表示违约概率。
        \item KMV模型使用历史违约率来确定预期违约频率，作为违约距离的函数。
    \end{itemize}
    
    \textbf{B选项错误。}
    \begin{itemize}
        \item 这个选项使用了股权和股权波动率，而不是资产和资产波动率。
        \item KMV模型关注的是资产价值，而不是股权价值。
    \end{itemize}
    
    \textbf{C选项错误。}
    \begin{itemize}
        \item 这个选项描述的是股价跌破阈值的概率，而不是违约距离的定义。
        \item 违约距离是一个标准化的距离指标，而不是概率。
    \end{itemize}
    
    \textbf{D选项错误。}
    \begin{itemize}
        \item 这个选项使用了杠杆和股价波动率的乘积，这不是违约距离的定义。
        \item 违约距离需要考虑资产价值和债务阈值，而不是杠杆和股价波动率。
    \end{itemize}
\end{explanation}

\checklist{违约距离的定义}

\begin{question}
    An analyst is using Moody's KMV model to estimate the distance to default of a large public firm, Sports Inc., a firm that designs, manufactures and sells athletic shoes. The firm's capital structure consists of USD 40 million in short-term debt, USD 20 million in long-term debt, and there are one million shares of stock currently trading at USD 10 per share. The asset volatility is 20\% per year. What is the normalized distance to default for Sports Inc.?
\end{question}

\begin{choices}
    \item 0.714
    \item 1.430
    \item 2.240
    \item 5.000
\end{choices}

\begin{correctanswer}
    B
\end{correctanswer}

\begin{explanation}
   
    
    \textbf{步骤1：确定参数}
    \begin{itemize}
        \item 短期负债：\$40 million
        \item 长期负债：\$20 million
        \item 股票数量：1 million股
        \item 每股价格：\$10
        \item 资产波动率：20\% = 0.20
    \end{itemize}

    \textbf{步骤2：计算债务阈值（K）}
    \begin{itemize}
        \item KMV模型中，债务阈值 = 短期负债 + 长期负债的一半
        \[
        K = 40 + 20 \times 0.5 = 40 + 10 = 50 \text{ million}
        \]
    \end{itemize}

    \textbf{步骤3：计算资产价值（A）}
    \begin{itemize}
        \item 股权市值 = 1{,}000{,}000 × \$10 = \$10 \text{ million}
        \item 总负债 = 40 + 20 = \$60 \text{ million}
        \item 资产价值：
        \[
        A = \text{股权市值} + \text{总负债} = 10 + 60 = 70 \text{ million}
        \]
    \end{itemize}

    \textbf{步骤4：计算资产波动率的绝对金额}
    \begin{itemize}
        \item 资产波动率的绝对金额：
        \[
        \sigma_A = A \times \sigma = 70 \times 0.20 = 14 \text{ million}
        \]
    \end{itemize}

    \textbf{步骤5：计算违约距离（DD）}
    \begin{itemize}
        \item 违约距离公式：
        \[
        DD = \frac{A - K}{\sigma_A}
        \]
        \item 代入数值：
        \[
        DD = \frac{70 - 50}{14} = \frac{20}{14} = 1.429 \approx 1.430
        \]
    \end{itemize}

\end{explanation}

\checklist{违约距离的计算}

\begin{question}
    David Wang is a senior analyst at an investment bank preparing research reports on the mining sector. Wang noted that one of the companies he follows recently increased risk to the firm's assets, which he expects will benefit equity holders to the detriment of debt holders. Which of the following concepts best describes the scenario in Wang's analysis?
\end{question}

\begin{choices}
    \item Coordination failures
    \item Adverse selection
    \item Risk shifting
    \item Principal-agent problem
\end{choices}

\begin{correctanswer}
    C
\end{correctanswer}

\begin{explanation}
    \textbf{A选项错误。}
    \begin{itemize}
        \item 协调失败（集体行动问题）发生在当一群人如果都采取某种行动会集体受益，但如果只有一个人采取该行动则不会受益的情况。本题描述的情况与此不同。
    \end{itemize}
    
    \textbf{B选项错误。}
    \begin{itemize}
        \item 逆向选择发生在交易双方信息不对称的情况下。本题描述的情况是风险转移，而不是信息不对称问题。
    \end{itemize}
    
    \textbf{C选项正确。}
    \begin{itemize}
        \item 风险转移（Risk shifting）发生在风险和收益从一组市场参与者转移到另一组持有不同公司资本结构头寸的参与者时。
        \item 在本例中，公司资产风险的增加会使股权持有人受益，而损害债务持有人的利益，因为债务持有人有固定收益但承担了更大的损失风险。
    \end{itemize}
    
    \textbf{D选项错误。}
    \begin{itemize}
        \item 委托-代理问题发生在委托人雇佣代理人执行特定职责，而代理人比委托人拥有更多信息的情况下。本题描述的情况是风险转移，而不是委托-代理问题。
    \end{itemize}
\end{explanation}

\checklist{风险转移的定义}


\section{考点2:信用评分与评级}
\setcounter{question}{0}

\begin{question}
    Michael Chen, FRM, is a rating agency analyst who is currently performing financial statement analysis on a major bank. Which of the following financial statements would be least useful for bank credit analysis?
\end{question}

\begin{choices}
    \item Balance sheet
    \item Income statement
    \item Statement of cash flows
    \item Statement of changes in capital funds
\end{choices}

\begin{correctanswer}
    C
\end{correctanswer}

\begin{explanation}
    \textbf{A选项错误。}
    \begin{itemize}
        \item 资产负债表对银行信用分析非常重要，因为它显示了银行的资产、负债和资本结构，这些都是评估银行信用风险的关键指标。
    \end{itemize}
    
    \textbf{B选项错误。}
    \begin{itemize}
        \item 利润表对银行信用分析也很重要，因为它显示了银行的收入、支出和盈利能力，这些都是评估银行偿债能力的重要指标。
    \end{itemize}
    
    \textbf{C选项正确。}
    \begin{itemize}
        \item 现金流量表在分析非金融实体时最有用，因为它可以区分经营、投资和融资活动的现金来源和用途。
        \item 然而，对于银行信用分析来说，现金流量表的用处较小，因为：
        \begin{enumerate}
            \item 银行的经营活动本身就是处理现金和现金等价物
            \item 银行的资产负债表和利润表已经包含了评估信用风险所需的大部分信息
            \item 银行的现金流模式与非金融实体有很大不同
        \end{enumerate}
    \end{itemize}
    
    \textbf{D选项错误。}
    \begin{itemize}
        \item 资本变动表对银行信用分析很重要，因为它显示了银行资本的变化情况，这是评估银行资本充足性和财务实力的重要指标。
    \end{itemize}
\end{explanation}

\checklist{现金流量表对银行信用分析的用处}

\begin{question}
    Each of the following items represents an example of qualitative information that would ideally be captured in assessing default probability except:
\end{question}

\begin{choices}
    \item Management's education and experience
    \item Internal controls associated with financial reporting
    \item Diversification of products and customers locally and globally
    \item Trends in throughput and other operational efficiency metrics
\end{choices}

\begin{correctanswer}
    D
\end{correctanswer}

\begin{explanation}
    \textbf{A选项错误。}
    \begin{itemize}
        \item 管理层的教育背景和经验是定性信息，难以用具体数字量化，需要主观评估和判断。
    \end{itemize}

    \textbf{B选项错误。}
    \begin{itemize}
        \item 与财务报告相关的内部控制是定性信息，虽然可用评估框架，但评估结果多为定性判断，需要主观评估。
    \end{itemize}

    \textbf{C选项错误。}
    \begin{itemize}
        \item 产品和客户在本地和全球的多元化是定性信息，虽然可以计算一些定量指标（如客户集中度），但多元化本身是定性概念，需要主观评估。
    \end{itemize}


    \textbf{D选项正确。}
    \begin{itemize}
        \item 吞吐量趋势和其他运营效率指标是定量信息，可以用具体数字衡量（如单位时间产量、设备利用率等），可以客观计算和比较，不依赖主观判断。
    \end{itemize}
\end{explanation}

\checklist{定性信息在信用分析中的应用}


\begin{question}
    Golin and Delhaise divide credit analysis into four areas according to borrower type:
    \begin{enumerate}
        \item Consumer credit analysis: evaluation of individual consumers' creditworthiness
        \item Corporate credit analysis: evaluation of nonfinancial companies
        \item Financial institution credit analysis: evaluation of financial companies
        \item Sovereign/municipal credit analysis: evaluation of credit risk of nations and governments
    \end{enumerate}
    According to Golin and Delhaise, each of the following is true about key features of credit analysis with respect to borrower type, except which is not true?
\end{question}

\begin{choices}
    \item Individuals (consumers): Credit analysis is amenable to automation and the use of scoring models and statistical tools to correlate risk to limited number of variables
    \item Non-financial corporations: Compared to consumers, tends to be more detailed and "hands-on" (i.e. less automated); key variables are likely to include liquidity, cash flow, near-term earnings capacity and profitability, solvency or capital position
    \item Financial Companies: In contrast to corporate (non-financial) credit analysis, qualitative analysis and asset quality are not important, but cash flow is a highly important (a "key indicator")
    \item Sovereigns: Includes analysis of country risk, which is primarily political dynamics and state of the economy; and systematic risk, which includes the regulatory regime and the financial system
\end{choices}

\begin{correctanswer}
    C
\end{correctanswer}

\begin{explanation}
   
    \textbf{A选项错误（陈述正确）。}
    \begin{itemize}
        \item 个人（消费者）信用分析适合自动化和使用评分模型，因为可以用有限的变量评估风险，可以使用统计工具建立标准化评分模型。
    \end{itemize}

    \textbf{B选项错误（陈述正确）。}
    \begin{itemize}
        \item 非金融企业的信用分析比消费者更详细和"手动"，需要评估流动性、现金流、短期和长期盈利能力、偿付能力和资本状况等。
    \end{itemize}

    \textbf{C选项正确（陈述错误）。}

    \begin{itemize}
        \item 陈述III错误：金融机构的信用分析中，定性分析和资产质量非常重要，而现金流不是关键指标。
        \item 资产质量对金融机构至关重要，直接影响信用风险；定性分析（管理质量、风险管理制度、监管合规性）同样重要。
        \item 现金流不是关键指标：金融机构的现金流模式与非金融企业不同，更关注盈利能力、资本充足率、流动性管理。
        \item 金融机构信用分析主要关注：盈利能力、资本充足率、资产质量、监管环境等。
    \end{itemize}


    \textbf{D选项错误（陈述正确）。}
    \begin{itemize}
        \item 主权信用分析包括国家风险分析（主要是政治动态和经济状况）和系统性风险分析（包括监管制度和金融体系）。
    \end{itemize}
\end{explanation}

\checklist{金融机构信用分析的关键要素}


\newpage


\section{考点3:违约可能性评估}
\setcounter{question}{0}
\begin{question}
    Which of the following statements is (are) not true regarding a credit default swap (CDS)?
    \begin{itemize}
        \item I.The purchaser of a CDS seeks credit protection
        \item II.Swap partners make payments based on a benchmark credit spread on the reference obligation
        \item III.Physical settlement provides the CDS buyer the par value of the reference obligation
        \item IV.A CDS has the same payoffs as a credit default put with installment payments
    \end{itemize}
\end{question}

\begin{choices}
    \item II only
    \item I and III
    \item III only
    \item II and IV
\end{choices}

\begin{correctanswer}
    A
\end{correctanswer}

\begin{explanation}
    \textbf{陈述I：正确。} 
    \begin{itemize}
        \item CDS的购买者确实寻求信用保护，通过购买CDS来对冲参考债务的信用风险。
    \end{itemize}
        
    \textbf{陈述II：错误。} 
    \begin{itemize}
        \item CDS的支付不是基于参考债务的基准信用利差。CDS买方在合约期限内或发生信用事件前，向CDS卖方支付固定费用（spread），这是预先确定的固定金额，通常以基点表示的年化费用，而不是基于基准信用利差。
    \end{itemize}
        
    \textbf{陈述III：正确。} 
    \begin{itemize}
        \item 实物结算时，CDS买方将参考债务交付给卖方，卖方支付面值给买方，因此买方获得了参考债务的面值。
    \end{itemize}
        
    \textbf{陈述IV：正确。} 
    \begin{itemize}
        \item CDS确实与分期付款的信用违约看跌期权有相同的收益结构，两者都提供信用保护。
    \end{itemize}
\end{explanation}

\checklist{CDS的定义和特点}

\begin{question}
    Which of the following statements is most likely correct about recovery rates?
\end{question}

\begin{choices}
    \item Recovery rates vary inversely with capital structure seniority
    \item Historical recovery rates are fairly constant across industries
    \item Recovery rates are highest during economic downturns
    \item Actual recovery rates may differ substantially from settled recovery rates
\end{choices}

\begin{correctanswer}
    D
\end{correctanswer}

\begin{explanation}
    \textbf{A选项错误。}
    \begin{itemize}
        \item 回收率与资本结构优先级成正比，而不是反比。优先级越高的债务，回收率越高；优先级越低的债务，回收率越低。
    \end{itemize}

    \textbf{B选项错误。}
    \begin{itemize}
        \item 历史回收率在不同行业之间存在显著差异。不同行业的资产特征、破产程序和经营周期敏感性不同，导致回收率差异较大。
    \end{itemize}

    \textbf{C选项错误。}
    \begin{itemize}
        \item 回收率在经济衰退期间最低，而不是最高。经济衰退时资产价值下降，破产程序可能更复杂，市场流动性可能更差，导致回收率下降。
    \end{itemize}

    \textbf{D选项正确。}
    \begin{itemize}
        \item 实际回收率可能与结算回收率有显著差异。结算回收发生在信用事件后不久（如CDS拍卖或违约债券出售），而实际回收可能发生在多年后，基于破产解决结果。破产程序可能持续很长时间，最终回收金额可能受到多种因素影响。
    \end{itemize}
\end{explanation}

\checklist{回收率的影响因素}


\begin{question}
    A 2-year credit default swap (CDS) specifying physical delivery defaults at the end of two years. If the reference asset is a \$200 million, 8.0\% ABC corporate bond, and the CDS spread is 125 basis points, the buyer of the CDS will:
\end{question}

\begin{choices}
    \item Receive payments of 800 basis points for the next two years
    \item Receive a payment of \$167.5 million
    \item Deliver the bond and receive a payment of \$200 million
    \item Continue to receive payments of 675 basis points for the next two years
\end{choices}

\begin{correctanswer}
    C
\end{correctanswer}

\begin{explanation}
   
    \textbf{A选项错误。}
    \begin{itemize}
        \item CDS买方不是接收支付，而是支付固定费用（CDS spread）给卖方。800个基点是债券的票面利率，不是CDS的支付。
    \end{itemize}

    \textbf{B选项错误。}
    \begin{itemize}
        \item \$167.5 million不是正确的支付金额。在实物交割的CDS中，买方会收到参考债务的面值（\$200 million），而不是\$167.5 million。
    \end{itemize}

    \textbf{C选项正确。}
    \begin{itemize}
        \item 在实物交割的CDS中，CDS买方将参考债务（ABC公司债券）交付给卖方，CDS买方收到参考债务的面值（\$200 million）。这是实物交割CDS的标准结算方式。
    \end{itemize}

    \textbf{D选项错误。}
    \begin{itemize}
        \item CDS买方不是接收支付，而是支付固定费用（CDS spread）给卖方。675个基点不是正确的CDS spread（实际是125个基点）。
    \end{itemize}
\end{explanation}

\checklist{CDS的结算方式}


\begin{question}
    Which of the following statements about credit default swaps is most accurate?
\end{question}

\begin{choices}
    \item CDSs transfer credit risk and market risk from the protection buyer to the protection seller
    \item CDSs transfer credit risk from the protection buyer to the issuer of the underlying credit
    \item Physical settlement requires knowledge of the post-default market price
    \item Cash settlement avoids the problem of a delivery squeeze
\end{choices}

\begin{correctanswer}
    D
\end{correctanswer}

\begin{explanation}
    \textbf{A选项错误。}
    \begin{itemize}
        \item CDS只转移信用风险，不转移市场风险。CDS主要针对信用事件提供保护，市场风险（如利率风险、汇率风险等）不在CDS的覆盖范围内。
    \end{itemize}

    \textbf{B选项错误。}
    \begin{itemize}
        \item CDS将信用风险从保护买方转移给保护卖方，而不是基础信用的发行人。保护卖方承担信用风险，基础信用的发行人仍然是信用风险的源头。
    \end{itemize}

    \textbf{C选项错误。}
    \begin{itemize}
        \item 实物结算不需要知道违约后的市场价格。实物结算基于参考债务的面值，市场价格与实物结算无关。
    \end{itemize}

    \textbf{D选项正确。}
    \begin{itemize}
        \item 现金结算避免了交割挤压问题。现金结算没有实际的证券交易，保护买方不需要在公开市场购买参考实体，因此不会因为大量购买导致价格上升。
    \end{itemize}
\end{explanation}

\checklist{现金结算的优点}


\begin{question}
    Economic capital calculations for credit risk assume a recovery rate (defined as 1-loss rate). Recovery rates are dependent on the business model of the underlying counterparty and its asset volatility in value and size. Under normal anticipated circumstances which of the following types of companies will have the highest recovery rate?
\end{question}

\begin{choices}
    \item An internet merchant of designer clothes
    \item A hedge fund
    \item An asset intensive manufacturing company
    \item A commodities trader
\end{choices}

\begin{correctanswer}
    C
\end{correctanswer}

\begin{explanation}
    \textbf{A选项错误。}
    \begin{itemize}
        \item 网上设计师服装零售商通常有形资产较少，主要资产是库存和应收账款，在违约时难以变现，因此回收率较低。
    \end{itemize}

    \textbf{B选项错误。}
    \begin{itemize}
        \item 对冲基金通常有形资产很少，主要资产是金融资产，资产价值波动较大，在违约时回收率较低。
    \end{itemize}

    \textbf{C选项正确。}
    \begin{itemize}
        \item 资产密集型制造公司通常拥有大量有形资产（如厂房、设备），资产价值相对稳定，在违约时容易估值和变现，因此回收率最高。例如，公用事业公司有很高的回收率，因为它们拥有大量有形资产，如发电厂。
    \end{itemize}

    \textbf{D选项错误。}
    \begin{itemize}
        \item 商品交易商通常有形资产较少，主要资产是商品库存，资产价值波动较大，在违约时回收率较低。
    \end{itemize}
\end{explanation}

\checklist{资产密集型制造公司的回收率优势}


\begin{question}
    An analyst has gathered the following information about ABC Inc. and DEF Inc. The respective credit ratings are AA and BBB with 1-year CDS spreads of 200 and 300 basis points each. The associated probabilities of default based on published reports are 10\% and 20\%, respectively. Which of the following statements about the recovery rates is most likely correct?
\end{question}

\begin{choices}
    \item The market implied recovery rates are equal
    \item The market implied recovery rate is higher for ABC
    \item The market implied recovery rate is lower for ABC
    \item The loss given default is higher for DEF
\end{choices}

\begin{correctanswer}
    C
\end{correctanswer}

\begin{explanation}
    \textbf{A选项错误。}
    \begin{itemize}
        \item 市场隐含回收率不相等：
        \begin{enumerate}
            \item ABC的回收率 = 80\%
            \item DEF的回收率 = 85\%
        \end{enumerate}
    \end{itemize}
    
    \textbf{B选项错误。}
    \begin{itemize}
        \item ABC的市场隐含回收率（80\%）低于DEF（85\%）
        \item 计算过程：
        \begin{enumerate}
            \item 信用利差 ≈ (1 - 回收率) × 违约概率
            \item ABC: 200 bps = (1 - RR) × 10\%, 所以RR = 80\%
            \item DEF: 300 bps = (1 - RR) × 20\%, 所以RR = 85\%
        \end{enumerate}
    \end{itemize}
    
    \textbf{C选项正确。}
    \begin{itemize}
        \item ABC的市场隐含回收率（80\%）确实低于DEF（85\%）
        \item 使用违约损失率（LGD）表示：
        \begin{enumerate}
            \item ABC的LGD = 20\%
            \item DEF的LGD = 15\%
        \end{enumerate}
    \end{itemize}
    
    \textbf{D选项错误。}
    \begin{itemize}
        \item DEF的违约损失率（15\%）低于ABC（20\%）
        \item 违约损失率 = 1 - 回收率
    \end{itemize}
\end{explanation}

\checklist{市场隐含回收率的计算方法}

\begin{question}
    The buyer of a credit-default swap (CDS):
\end{question}

\begin{choices}
    \item Has no risk exposure to the combined holdings of the CDS and reference obligation
    \item Is subject to the credit risk of the seller
    \item Must place collateral with the seller of the CDS to ensure performance
    \item Can surrender the CDS for the value of the accumulated payments
\end{choices}

\begin{correctanswer}
    B
\end{correctanswer}

\begin{explanation}
    \textbf{A选项错误。}
    \begin{itemize}
        \item CDS买方仍然面临风险，包括参考债务的信用风险、CDS卖方的信用风险（对手方风险）和基础风险（basis risk）。
    \end{itemize}

    \textbf{B选项正确。}
    \begin{itemize}
        \item CDS买方确实面临卖方的信用风险（对手方风险）。如果参考债务违约，卖方可能无法支付所需款项。两个因素影响这种风险：参考债务违约概率随时间增加，卖方违约概率也可能增加；经济和运营因素可能导致参考债务违约和卖方违约之间的相关性增加。
    \end{itemize}

    \textbf{C选项错误。}
    \begin{itemize}
        \item CDS买方不需要向卖方提供抵押品。买方支付固定费用（CDS spread），抵押品要求通常由卖方提供。
    \end{itemize}

    \textbf{D选项错误。}
    \begin{itemize}
        \item CDS买方不能以累积支付的价值放弃CDS。CDS是双边合约，没有提前终止的自动机制，需要双方同意才能终止。
    \end{itemize}

\end{explanation}

\checklist{CDS买方的风险暴露}

\begin{question}
    With regard to the factors that determine recovery rates of traded bonds, which of the following statements is false?
\end{question}

\begin{choices}
    \item The liquidity of the debtor's collateral has an impact of recovery rates of the firm's traded bonds
    \item The presence of more bank loans increases the recovery rates of the firm's traded bonds
    \item The seniority of debtor's debt has an impact of recovery rates of the firm's traded bonds
    \item The value of the debtor's collateral has an impact of recovery rates of the firm's traded bonds
\end{choices}

\begin{correctanswer}
    B
\end{correctanswer}

\begin{explanation}
    \textbf{A选项错误（陈述正确）。}
    \begin{itemize}
        \item 债务人抵押品的流动性确实会影响交易债券的回收率。流动性高的抵押品更容易变现，流动性低的抵押品可能难以在违约时快速出售。
    \end{itemize}

    \textbf{B选项正确（陈述错误）。}
    \begin{itemize}
        \item 更多的银行贷款实际上会降低交易债券的回收率，而不是提高。银行贷款通常是优先级债务，有最高级别的担保。公司使用的银行融资越多，交易债券的回收率就越低，因为银行债务优先于交易债券获得偿付。
    \end{itemize}

    \textbf{C选项错误（陈述正确）。}
    \begin{itemize}
        \item 债务人债务的优先级确实会影响交易债券的回收率。优先级定义了公司违约时债权人对资产的优先权顺序，优先级高的债务通常有更高的回收率。
    \end{itemize}

    \textbf{D选项错误（陈述正确）。}
    \begin{itemize}
        \item 债务人抵押品的价值确实会影响交易债券的回收率。抵押品价值越高，回收率可能越高，抵押品价值是回收率的重要决定因素。
    \end{itemize}
\end{explanation}

\checklist{影响交易债券回收率的因素}


\begin{question}
    A six-year CDS on an AA-rated issuer is offered at 150bp with semiannual payments while the yield on a six-year semiannual coupon bond of this issuer is 8\%. There is no counterparty risk on the CDS. The annualized LIBOR rate paid every six months is 4.6\% for all maturities. Which strategy would exploit the arbitrage opportunity? How much would your return exceed LIBOR?
\end{question}

\begin{choices}
    \item Buy the bond and the CDS with a risk-free gain of 1.9\%
    \item Buy the bond and the CDS with a risk-free gain of 0.32\%
    \item Short the bond and sell CDS protection with a risk-free gain of 4.97\%
    \item There is no arbitrage opportunity as any apparent risk-free profit is necessarily compensation for being exposed to the credit risk of the issuer
\end{choices}

\begin{correctanswer}
    A
\end{correctanswer}

\begin{explanation}

    \textbf{A选项正确。}

    \begin{itemize}
        \item 步骤1：计算债券的信用利差
        \begin{itemize}
            \item 债券收益率：8\% = 800 bp
            \item LIBOR（无风险利率）：4.6\% = 460 bp
            \item 信用利差：800 - 460 = 340 bp
        \end{itemize}

        \item 步骤2：确定套利策略
        \begin{itemize}
            \item 买入债券并获得8\%的收益率
            \item 买入CDS保护，支付150 bp的费用来对冲信用风险
            \item 由于CDS没有交易对手风险，这个策略是无风险的
        \end{itemize}

        \item 步骤3：计算超额收益
        \begin{itemize}
            \item 净收益 = 信用利差 - CDS费用 = 340 bp - 150 bp = 190 bp = 1.9\%
            \item 这个收益是相对于LIBOR的超额收益，且是无风险的
        \end{itemize}
    \end{itemize}

    \textbf{B选项错误。}
    \begin{itemize}
        \item 0.32\%的收益计算错误，正确的收益应该是1.9\%。
    \end{itemize}

    \textbf{C选项错误。}
    \begin{itemize}
        \item 做空债券和出售CDS保护会产生损失，而不是收益。4.97\%的收益计算错误。
    \end{itemize}

    \textbf{D选项错误。}
    \begin{itemize}
        \item 确实存在套利机会。由于CDS没有交易对手风险，这个套利是真正无风险的，不是对信用风险的补偿。
    \end{itemize}

\end{explanation}

\checklist{CDS套利策略}

\begin{question}
    The Merton model and the Moody's KMV model use different approaches to determine the probability of default. Which of the following is consistent with Moody's KMV model?
\end{question}

\begin{choices}
    \item The distance to default is 1.96, so there is a 2.5\% probability of default
    \item The distance to default is 1.96, so there is a 5.0\% probability of default
    \item The historical frequency of default for corporate bonds has been 6\%. Updating this with Altman's Z-score analysis would provide a probability of default that is somewhat different than 6\%
    \item The distance to default is 1.96 and, historically, 1.2\% of firms with this characterization have defaulted, so there is a 1.2\% probability of default
\end{choices}

\begin{correctanswer}
    D
\end{correctanswer}

\begin{explanation}
    \textbf{A选项错误。}
    \begin{itemize}
        \item 2.5\%的违约概率是基于正态分布假设（距离违约1.96对应正态分布的单尾2.5\%概率）。Moody's KMV模型不使用正态分布假设，而是使用历史违约数据来确定违约概率。
    \end{itemize}

    \textbf{B选项错误。}
    \begin{itemize}
        \item 5.0\%的违约概率也是基于正态分布假设（双尾5\%）。这与Moody's KMV模型的方法不符，模型使用历史违约频率而不是理论概率。
    \end{itemize}

    \textbf{C选项错误。}
    \begin{itemize}
        \item 使用Altman's Z-score分析不是Moody's KMV模型的方法。模型使用距离违约（distance to default）作为主要指标，不依赖于Z-score分析。
    \end{itemize}

    \textbf{D选项正确。}
    \begin{itemize}
        \item Moody's KMV模型评估具有相似违约距离的公司的历史违约频率，并使用这个历史频率作为违约概率，不依赖于理论分布假设。
        \item 在这个例子中，违约距离为1.96，历史上具有这个特征的公司有1.2\%违约，因此违约概率为1.2\%。
    \end{itemize}
\end{explanation}

\checklist{Moody's KMV模型的违约概率确定方法}


\begin{question}
    A corporate bond with a face value of \$1,000 has a remaining maturity of 10 years. Using the Merton model, the current value of the bond is calculated at \$750. Assuming that the risk-free rate is equal to 3\%, what is the credit spread for this bond?
\end{question}

\begin{choices}
    \item 35 bps
    \item 42 bps
    \item 58 bps
    \item 65 bps
\end{choices}

\begin{correctanswer}
    C
\end{correctanswer}

\begin{explanation}
    

    
    \textbf{C选项正确。}
    \begin{itemize}
        \item 使用Merton模型的信用利差公式：
        \[
        \begin{aligned}
        \text{信用利差} &= -\left(\frac{1}{T-t}\right) \ln\left(\frac{D}{F}\right) - r \\
        &= -\left(\frac{1}{10}\right) \times \ln\left(\frac{750}{1000}\right) - 0.03 \\
        &= 58\ \text{个基点}
        \end{aligned}
        \]
        \item 其中：
        \begin{enumerate}
            \item $T-t = 10$年（剩余期限）
            \item $D = \$750$（债券现值）
            \item $F = \$1,000$（面值）
            \item $r = 3\%$（无风险利率）
        \end{enumerate}
    \end{itemize}
    
   
\end{explanation}

\checklist{Merton模型的信用利差计算}

\begin{question}
    The capital structure of ABC Corporation consists of two parts, one 4 year bond with face value of \$120 million and the rest is equity. The current market value of the firm assets is \$150 and the expected rate of change of the firm's value is 20\%. The volatility is 25\%. The firm's risk management division estimates the distance to default using Merton model. Given the distance to default, the estimated physical default probability is? (N(1.85) = 96.78\%; N(2.45) = 99.29\%)
\end{question}

\begin{choices}
    \item 3.22\%
    \item 15.36\%
    \item 15.38\%
    \item 28.45\%
\end{choices}

\begin{correctanswer}
    A
\end{correctanswer}

\begin{explanation}
    \textbf{A选项正确。}
    \begin{itemize}
        \item 使用Merton模型计算违约距离：
        \[
        \begin{aligned}
        d_2 &= \frac{\ln\left[\frac{150}{120} \times e^{-20\% \times 4}\right]}{0.25 \times \sqrt{4}} - \frac{1}{2} \times 0.25 \times \sqrt{4} \\
        &= 1.85
        \end{aligned}
        \]
        \item 实际违约概率 = $N(-d_2)$ = 3.22\%
        \item 其中：
        \begin{enumerate}
            \item 公司资产价值 = \$150 million
            \item 债券面值 = \$120 million
            \item 预期增长率 = 20\%
            \item 波动率 = 25\%
            \item 期限 = 4年
        \end{enumerate}
    \end{itemize}

\end{explanation}

\checklist{使用Merton模型计算违约距离}


\begin{question}
    Which of the following statements regarding the Merton model is true?
\end{question}

\begin{choices}
    \item A firm with multiple debt issues that mature at different times is easy to value with the Merton model
    \item The Merton model assumes a lognormal distribution and constant variance for changes in firm value
    \item The Merton model is able to predict default because it allows for default surprises (i.e., jumps)
    \item Empirical results indicate that the Merton model is able to predict default better than naive models for investment grade bonds
\end{choices}

\begin{correctanswer}
    B
\end{correctanswer}

\begin{explanation}
    \textbf{A选项错误。}
    \begin{itemize}
        \item Merton模型难以处理多种债务工具。大多数公司有不同到期日的债务，债务可能有不同的票面利率，但模型假设单一零息债券，这与现实不符。
    \end{itemize}

    \textbf{B选项正确。}
    \begin{itemize}
        \item Merton模型的基本假设是公司价值服从对数正态分布，且公司价值变化具有恒定方差。这些假设是模型的核心特征。
    \end{itemize}

    \textbf{C选项错误。}
    \begin{itemize}
        \item Merton模型不允许公司价值跳跃。模型假设公司价值连续变化，不能处理突然的违约事件，这使得违约预测过于可预测。
    \end{itemize}

    \textbf{D选项错误。}
    \begin{itemize}
        \item 实证研究结果表明，对于投资级债券，简单模型（假设债务无风险）比Merton模型更有效。Merton模型对非投资级债券的违约预测更有效，这与选项描述相反。
    \end{itemize}
\end{explanation}

\checklist{Merton模型的基本假设}





\begin{question}
    A default swap acts like a:
\end{question}

\begin{choices}
    \item Call option on the reference obligation for the buyer of the swap
    \item Put option on the reference obligation for the buyer of the swap
    \item Look back option on the reference obligation for the buyer of the swap
    \item Up-and-out option on the reference obligation for the buyer of the swap
\end{choices}

\begin{correctanswer}
    B
\end{correctanswer}

\begin{explanation}
    \textbf{A选项错误。}
    \begin{itemize}
        \item 看涨期权（Call Option）的特征是买方有权以固定价格买入资产，当资产价格上升时获利，这与违约互换的功能不符。
    \end{itemize}

    \textbf{B选项正确。}
    \begin{itemize}
        \item 违约互换类似于看跌期权（Put Option）。买方在违约时获得支付，限制了买方的下行风险，类似于看跌期权的保护功能。当参考债务价值下降时提供保护，在违约事件发生时获得补偿，类似于看跌期权的行权机制。
    \end{itemize}

    \textbf{C选项错误。}
    \begin{itemize}
        \item 回顾期权（Look Back Option）的特征是基于资产价格的历史表现，与违约事件无关，不适用于违约互换。
    \end{itemize}

    \textbf{D选项错误。}
    \begin{itemize}
        \item 向上敲出期权（Up-and-out Option）的特征是当资产价格超过特定水平时失效，与违约事件无关，不适用于违约互换。
    \end{itemize}
\end{explanation}

\checklist{违约互换的性质}



\begin{question}
    Which of the following statements is most likely correct about recovery rates?
\end{question}

\begin{choices}
    \item Recovery rates are higher for secured debt than for unsecured debt
    \item Recovery rates remain stable across different economic cycles
    \item Recovery rates are higher for subordinated debt than for senior debt
    \item Actual recovery rates may differ substantially from settled recovery rates
\end{choices}

\begin{correctanswer}
    D
\end{correctanswer}

\begin{explanation}
    \textbf{A选项错误。}
    \begin{itemize}
        \item 虽然"有担保债务的回收率高于无担保债务"这一陈述通常是正确的，但这不是题目要求的"最准确"的描述。题目要求找出关于回收率最可能正确的陈述，而D选项提供了更重要的洞察。【疑惑】
    \end{itemize}

    \textbf{B选项错误。}
    \begin{itemize}
        \item 回收率在不同经济周期中并不稳定。经济衰退期间，资产价值下降、市场流动性差，回收率通常较低；经济繁荣期间回收率通常较高。
    \end{itemize}

    \textbf{C选项错误。}
    \begin{itemize}
        \item 次级债务的回收率低于优先债务，而不是更高。在违约清算时，优先债务优先获得偿付，因此回收率更高。
    \end{itemize}

    \textbf{D选项正确。}
    \begin{itemize}
        \item 实际回收率确实可能与结算回收率有显著差异。结算回收率发生在信用事件后不久（如CDS拍卖或违约债券的快速出售），而实际回收率可能发生在多年后的破产程序结束时，基于最终的资产清算结果，两者可能存在显著差异。
    \end{itemize}
\end{explanation}

\checklist{回收率的影响因素}

\begin{question}
    Suppose a firm has two debt issues outstanding. One is a senior debt issue that matures in four years with a principal amount of \$120 million. The other is a subordinate debt issue that also matures in four years with a principal amount of \$60 million. The annual interest rate is 4\% and the volatility of the firm value is estimated to be 20\%. If the volatility of the firm value declines in the Merton model then which of the following statements is true?
\end{question}

\begin{choices}
    \item If the firm is experiencing financial distress (low firm value), then the value of senior debt will increase while the values of subordinate debt and equity will both decline
    \item If the firm is not experiencing financial distress (high firm value), then the value of senior debt and subordinate debt and equity will increase
    \item If the firm is experiencing financial distress (low firm value), then the value of senior debt and subordinate debt will increase while equity values will decline
    \item If the firm is not experiencing financial distress (high firm value), then the value of senior debt will increase while the values of subordinate debt and equity will both decline
\end{choices}

\begin{correctanswer}
    A
\end{correctanswer}

\begin{explanation}

    \textbf{A选项正确。}
    \begin{itemize}
        \item 困境情形（公司价值低）：股权视作对资产的看涨期权，波动率下降会降低期权价值→股权下降；次级债更“股权化”→价值下降；高级债受益于违约风险下降→价值上升。
    \end{itemize}

    \textbf{B选项错误。}
    \begin{itemize}
        \item 非困境情形（公司价值高）：两类债务的违约风险均下降→高级债与次级债价值均上升；但股权的期权性价值随波动率下降而下降，不会与两类债务同向上升。
    \end{itemize}

    \textbf{C选项错误。}
    \begin{itemize}
        \item 困境情形下，次级债更接近股权性质，波动率下降使其价值下降，不会与高级债同涨。
    \end{itemize}

    \textbf{D选项错误。}
    \begin{itemize}
        \item 非困境情形下，两类债务价值上升，不存在“高级债上升而次级债与股权同时下降”的组合关系。
    \end{itemize}

\end{explanation}

\checklist{Merton模型的应用}

\begin{question}
    Which of the following statements regarding the Merton model is true?
\end{question}

\begin{choices}
    \item A firm with multiple debt issues that mature at different times is easy to value with the Merton model
    \item The Merton model assumes a lognormal distribution and constant variance for changes in firm value
    \item The Merton model is able to predict default because it allows for default surprises (i.e., jumps)
    \item Empirical results indicate that the Merton model is able to predict default better than naive models for investment grade bonds
\end{choices}

\begin{correctanswer}
    B
\end{correctanswer}

\begin{explanation}

    \textbf{A选项错误。}
    \begin{itemize}
        \item 模型假设资本结构可等价为“股权 + 单一到期的零息债”；现实中多数公司存在多期、多票息债务与嵌入条款，直接套用估值困难。
    \end{itemize}

    \textbf{B选项正确。}
    \begin{itemize}
        \item 核心假设：公司总资产遵循几何布朗运动（对数正态），波动率恒定；违约仅在到期点因资产低于债务面值而触发。
    \end{itemize}

    \textbf{C选项错误。}
    \begin{itemize}
        \item 不含跳跃项，资产路径连续；模型不允许“违约惊讶（跳跃）”，因此难以刻画突发性违约。
    \end{itemize}

    \textbf{D选项错误。}
    \begin{itemize}
        \item 证据表明：对投资级主体，简单经验模型常优于Merton；Merton对高风险（非投级）违约区分度更有优势。
    \end{itemize}

\end{explanation}

\checklist{Merton 模型的核心假设}

\begin{question}
    Each of the following items represents an example of qualitative information that would ideally be captured in assessing default probability except:
\end{question}

\begin{choices}
    \item Management's education and experience
    \item Internal controls associated with financial reporting
    \item Diversification of products and customers locally and globally
    \item Trends in throughput and other operational efficiency metrics
\end{choices}

\begin{correctanswer}
    D
\end{correctanswer}

\begin{explanation}

    \textbf{A选项错误。}
    \begin{itemize}
        \item 管理层的教育与经验属于定性信息，难以用统一数值刻度准确量化，评估依赖主观判断，但对违约概率具有重要影响。
    \end{itemize}

    \textbf{B选项错误。}
    \begin{itemize}
        \item 与财务报告相关的内部控制属于定性信息，需基于框架与专业判断评估有效性，难以用单一数值衡量，但直接影响财务信息可靠性与违约风险。
    \end{itemize}

    \textbf{C选项错误。}
    \begin{itemize}
        \item 产品与客户在本地与全球的多元化属于定性信息，其对集中度与韧性的影响需综合评估，通常无法用单一量表精确度量。
    \end{itemize}

    \textbf{D选项正确。}
    \begin{itemize}
        \item 吞吐量趋势与其他运营效率指标属于定量信息，可用客观数值（产能利用率、单位成本、周转率等）持续跟踪、统计检验并横向比较，因此不属于定性信息。
    \end{itemize}

\end{explanation}

\checklist{定性信息与定量信息的区别}


\begin{question}
    The rate parameter in the exponential distribution measures the rate at which it takes an event to occur. In the context of waiting for a company to default, the rate parameter is known as the hazard rate and indicates the rate at which default will arrive. Which of the following statements about hazard rates (i.e., default intensity) is correct assuming a constant default intensity of 0.15?
\end{question}

\begin{choices}
    \item The cumulative default probability after 3 periods = $e^{-(0.15)3}$
    \item The conditional default probability after 1 period = $1 - e^{-(0.15)}$
    \item The unconditional default probability is memoryless
    \item The conditional default probability is memoryless
\end{choices}

\begin{correctanswer}
    D
\end{correctanswer}

\begin{explanation}
    \textbf{A选项错误。}
    \begin{itemize}
        \item 指数分布的生存函数为$S(t)=e^{-\lambda t}$，累计违约概率为$F(t)=1-e^{-\lambda t}$。
        \item 题干给式$e^{-(0.15)\cdot 3}$是生存概率$S(3)$，而非累计违约概率$F(3)$；正确应为$1-e^{-0.45}$。
    \end{itemize}

    \textbf{B选项错误。}
    \begin{itemize}
        \item “条件违约概率 after 1 period”表述不严谨。若指给定已存活到某时点，未来1期的条件违约概率确为$1-e^{-0.15}$；但若指从0到1期的“after 1 period”累计违约，属无条件违约概率，与“条件”不符。
    \end{itemize}

    \textbf{C选项错误。}
    \begin{itemize}
        \item “无条件违约概率无记忆”错误。指数分布的无记忆性适用于条件概率：$P(T>t+s\mid T>s)=P(T>t)$，而非无条件$F(t)$。
    \end{itemize}

    \textbf{D选项正确。}
    \begin{itemize}
        \item 指数分布（恒定违约强度$\lambda=0.15$）具有无记忆性：已存活至$s$后，未来$\Delta$期的条件违约概率为$1-e^{-\lambda \Delta}$，仅与$\Delta$有关，与过去$s$无关。
    \end{itemize}
\end{explanation}

\checklist{指数分布的无记忆性}


\begin{question}
    Using the Merton model, calculate the current value of a firm's equity and debt given that the current value of the firm is \$120 million, the principal amount due in four years on the zero-coupon bond is \$100 million, the annual interest rate is 8\%, and the volatility of the firm is 25\%.
\end{question}

\begin{choices}
    \item \$100 million in debt and \$20 million in equity
    \item \$65.32 million in debt and \$54.68 million in equity
    \item \$62.45 million in debt and \$57.55 million in equity
    \item \$45.28 million in debt and \$74.72 million in equity
\end{choices}

\begin{correctanswer}
    C
\end{correctanswer}

\begin{explanation}
   
    \begin{itemize}
        \item 参数与折现债务：
        \[
        V_0=120,\quad F=100,\quad r=8\%,\quad \sigma=25\%,\quad T=4,\quad F e^{-rT}=100e^{-0.08\times4}=73.50
        \]
        \item Merton公式（股权视作看涨期权）：
        \[
        S_0=V_0 N(d_1)-F e^{-rT} N(d_2),\quad d_1=\frac{\ln\!\left(\frac{V_0}{F e^{-rT}}\right)+\tfrac{1}{2}\sigma^2 T}{\sigma\sqrt{T}},\quad d_2=d_1-\sigma\sqrt{T}
        \]
        \item 代入计算：
        \[
        d_1=\frac{\ln(120/73.50)+0.5\times0.25^2\times4}{0.25\sqrt{4}}=1.2456,\quad d_2=1.2456-0.5=0.7456
        \]
        \[
        N(d_1)\approx0.8934,\quad N(d_2)\approx0.6915
        \]
        \[
        S_0=120\times0.8934-73.50\times0.6915=107.21-50.83=57.55
        \]
        \item 债务价值（资产减股权）：
        \[
        D_0=V_0-S_0=120-57.55=62.45
        \]
        \item 结论：债务约\$62.45m，股权约\$57.55m（对应C）。
    \end{itemize}
    
 
\end{explanation}

\checklist{使用Merton模型计算股权和债务的价值}

\begin{question}
    A financial analyst wants to estimate the value of a company's equity using the Merton model. The company pays no dividends and has a market value of assets equal to EUR 80 million. Its only debt, a zero-coupon bond maturing in 3 years, has a face value of EUR 50 million and a current price of EUR 62.50 per EUR 100 of face value. Additional information is given below:
    • N(d1): 0.9625
    • N(d2): 0.9012
    • Annual risk-free interest rate: 4\%
    What is the value of the company's equity?
\end{question}

\begin{choices}
    \item EUR 28.6 million
    \item EUR 32.4 million
    \item EUR 42.8 million
    \item EUR 47.5 million
\end{choices}

\begin{correctanswer}
    D
\end{correctanswer}

\begin{explanation}
    \textbf{D选项正确。}
    \begin{itemize}
        \item 使用Merton模型计算：
        \[
        \begin{aligned}
        S_t &= V_tN(d_1) - Ke^{-r(T-t)}N(d_2) \\
        &= 80 \times 0.9625 - 50 \times e^{-0.04 \times 3} \times 0.9012 \\
        &= 77 - 29.5 = 47.5
        \end{aligned}
        \]
    \end{itemize}
\end{explanation}

\checklist{使用Merton模型计算股权和债务的价值}






\begin{question}
    A large regional bank is changing the method it uses to calculate credit risk capital from the standardized approach to the Foundation Internal Ratings-Based (IRB) approach. A risk analyst is asked to summarize the process of calculating the capital charge under the new approach. Which of the following actions would the bank be required to take as part of the Foundation IRB approach?
\end{question}

\begin{choices}
    \item Apply an adjustment factor that reduces the risk weight for credit exposures with maturities longer than 1 year
    \item Develop a matrix to estimate the correlations between different classes of credit exposures
    \item Assume that correlations between high-risk corporate bonds are constant across all market conditions
    \item Provide an internally modeled default probability for each credit exposure
\end{choices}

\begin{correctanswer}
    D
\end{correctanswer}

\begin{explanation}

    \textbf{A选项错误。}
    \begin{itemize}
        \item 期限调整因子不是基础IRB方法的要求，而是高级IRB方法的特征。基础IRB方法使用监管机构提供的标准期限调整。
    \end{itemize}

    \textbf{B选项错误。}
    \begin{itemize}
        \item 相关性矩阵不是基础IRB方法的要求，而是高级IRB方法的特征。基础IRB方法使用监管机构提供的标准相关性参数。
    \end{itemize}

    \textbf{C选项错误。}
    \begin{itemize}
        \item 假设相关性恒定不是基础IRB方法的要求。相关性会随市场条件变化，基础IRB方法使用监管机构提供的标准相关性参数，这些参数可能反映不同市场条件下的相关性。
    \end{itemize}

    \textbf{D选项正确。}
    \begin{itemize}
        \item 基础IRB方法的核心要求是必须为每个信用敞口提供内部建模的违约概率（PD）。其他参数（如违约损失率LGD、违约风险敞口EAD、期限M）使用监管机构提供的标准值。
    \end{itemize}

\end{explanation}

\checklist{基础IRB方法的要求}   


\begin{question}
    A credit analyst is estimating a bank's economic capital for credit risk at the 99.9\% confidence level. The bank reports that its current credit risk portfolio has the following parameters:
    \begin{itemize}
    \item Total exposure amount: GBP 250 million
    \item Probability of default (PD): 2.5\%
    \item Loss rate (LR): 12\%
    \item Standard deviation of the PD: 8\%
    \item Standard deviation of the LR: 15\%
    \item Capital multiplier for the portfolio at the 99.9\% confidence level: 7.50
    \end{itemize}
    What amount of economic capital should the bank hold for its credit risk portfolio?
\end{question}

\begin{choices}
    \item GBP 5.25 million
    \item GBP 6.75 million
    \item GBP 37.50 million
    \item GBP 45.00 million
\end{choices}

\begin{correctanswer}
    D
\end{correctanswer}

\begin{explanation}
    
    \textbf{D选项正确。}
    \begin{itemize}
        \item 经济资本计算：
        \[
        \begin{aligned}
        \text{经济资本} &= \text{敞口金额} \times \text{资本乘数} \times \sqrt{\text{PD标准差}^2 + \text{LR标准差}^2} \\
        &= 250 \times 7.50 \times \sqrt{0.08^2 + 0.15^2} \\
        &= 250 \times 7.50 \times 0.17 \\
        &= 45.00
        \end{aligned}
        \]
    \end{itemize}
\end{explanation}

\checklist{经济资本的计算}


\begin{question}
    A team of credit analysts is using the Merton model to estimate the credit rating of a non-dividend-paying firm. The team compiles relevant information on the firm shown below:
    \begin{itemize}
    \item Value of total assets: AUD 250 million
    \item Long-term debt: AUD 120 million
    \item Short-term debt: AUD 0
    \item Expected return per year on firm assets: 12\%
    \item Asset volatility per year: 25\%
\end{itemize}
    The analysts use the following S\&P rating schedule provided by the investment bank to draw their conclusion:

    \begin{table}[H]
        \centering
        \renewcommand{\arraystretch}{1.2}
        \begin{tabular}{|c|c|}
            \hline
            \textbf{Range of distance to default} & \textbf{S\&P rating class} \\
            \hline
            10.0 and greater & AAA \\
            \hline
            8.0 - 9.9 & AA \\
            \hline
            6.0 - 7.9 & A \\
            \hline
            4.0 - 5.9 & BBB \\
            \hline
            2.0 - 3.9 & BB \\
            \hline
            1.0 - 1.9 & B \\
            \hline
            0.5 - 0.9 & C \\
            \hline
        \end{tabular}
        \caption{距离违约区间与S\&P评级对应表}
    \end{table}

    What is the suggested credit rating of the firm at a 1-year horizon using the S\&P rating schedule?
\end{question}

\begin{choices}
    \item A
    \item BB
    \item B
    \item C
\end{choices}

\begin{correctanswer}
    B
\end{correctanswer}

\begin{explanation}
   
    \textbf{B选项正确。}
    \begin{itemize}
        \item 违约距离计算：
        \[
        \begin{aligned}
        \text{违约距离} &= \frac{\ln\left(\frac{250}{120}\right) + (0.12 - 0.5 \times 0.25^2) \times 1}{0.25 \times \sqrt{1}} \\
        &= \frac{0.733 + 0.089}{0.25} \\
        &= 3.29
        \end{aligned}
        \]
        \item 根据S\&P评级表：
        \begin{itemize}
            \item 违约距离 = 3.29
            \item 落在2.0 - 3.9区间
            \item 对应评级为BB
        \end{itemize}
    \end{itemize}
    
 
\end{explanation}

\checklist{使用Merton模型计算信用评级}





\section{考点4:信用评分和零售信用风险管理}
\setcounter{question}{0}
\begin{question}
    Smith and Johnson divide credit analysis into four areas according to borrower type:
    Ⅰ.Consumer credit analysis is the evaluation of the creditworthiness of individual consumers;
    Ⅱ.Corporate credit analysis is the evaluation of nonfinancial companies such as manufacturers,
    and nonfinancial service providers; Ⅲ.Financial institution credit analysis is the evaluation
    of financial companies including banks and nonbank financial institutions, such as insurance
    companies and investment funds; Ⅳ.Sovereign/municipal credit analysis is the evaluation of the
    credit risk associated with the financial obligations of nations, subnational governments, and
    public authorities, as well as the impact of such risks on obligations of nonstate entities
    operating in specific jurisdictions. According to Smith and Johnson, each of the following is
    true about key features of credit analysis with respect to borrower type, except which is not
    true?
\end{question}

\begin{choices}
    \item Individuals (consumers): Credit analysis is amenable to automation and the use of scoring
    models and statistical tools to correlate risk to limited number of variables
    \item Non-financial corporations: Compared to consumers, tends to be more detailed and "hands-on"
    (i.e. less automated); key variables are likely to include liquidity, cash flow, near-term
    earnings capacity and profitability, solvency or capital position
    \item Financial Companies: In contrast to corporate (non-financial) credit analysis, qualitative
    analysis and asset quality are not important, but cash flow is a highly important (a "key
    indicator")
    \item Sovereigns: Includes analysis of country risk, which is primarily political dynamics and
    state of the economy; and systematic risk, which includes the regulatory regime and the
    financial system
\end{choices}

\begin{correctanswer}
    C
\end{correctanswer}

\begin{explanation}
    \textbf{A选项正确。}
    \begin{itemize}
    \item 消费者信用分析适合自动化和使用评分模型，可以用统计工具将风险与有限数量的变量相关联。这是正确的，因为消费者信用分析通常涉及大量标准化数据，适合使用自动化工具和统计模型。
    \end{itemize}

    \textbf{B选项正确。}
    \begin{itemize}
    \item 非金融企业信用分析比消费者分析更详细且需要更多人工参与，关键变量包括流动性、现金流、盈利能力和资本状况。这是正确的，因为企业分析需要更全面的评估和更多定性因素。
    \end{itemize}

    \textbf{C选项错误。}
    \begin{itemize}
    \item 金融机构信用分析中，资产质量非常重要，而现金流不是关键指标。这与选项描述相反，因为金融机构的资产质量是评估其信用风险的关键因素，而现金流分析在金融机构信用评估中相对不那么重要。
    \end{itemize}

    \textbf{D选项正确。}
    \begin{itemize}
    \item 主权国家信用分析包括国家风险分析（政治动态和经济状况）和系统性风险分析（监管制度和金融体系）。这是正确的，因为主权信用分析需要同时考虑政治、经济和金融体系等多个维度。
    \end{itemize}
\end{explanation}

\checklist{不同借款人类型的信用分析特点}


\begin{question}
    Which of the following is not an example or element of predatory lending:
\end{question}

\begin{choices}
    \item Lender makes unaffordable loans based on borrower assets rather than ability to repay
    \item Lender induces borrower to repeatedly refinance in order to collect fees and charge high points
    \item Borrower misrepresents income or employment in mortgage application
    \item Lender engages in deception to conceal true nature of loan; e.g., deceives borrower into thinking loan is fixed-rate when mortgage is actually an adjustable-rate
\end{choices}

\begin{correctanswer}
    C
\end{correctanswer}

\begin{explanation}
    \textbf{A选项错误。}
    \begin{itemize}
    \item 这是掠夺性贷款的一个典型特征。贷款人基于借款人的资产而不是还款能力发放贷款，这可能导致借款人无法负担贷款，属于掠夺性贷款行为。
    \end{itemize}

    \textbf{B选项错误。}
    \begin{itemize}
    \item 这是掠夺性贷款的另一个特征。贷款人诱导借款人反复再融资以收取费用和高额点数，这种行为被称为"贷款翻新"（loan flipping），是掠夺性贷款的常见手段。
    \end{itemize}

    \textbf{C选项正确。}
    \begin{itemize}
    \item 这是掠夺性借款（predatory borrowing）的例子，而不是掠夺性贷款（predatory lending）。借款人通过伪造收入或就业信息来获得贷款，这是借款人的欺诈行为，而不是贷款人的掠夺性行为。
    \end{itemize}

    \textbf{D选项错误。}
    \begin{itemize}
    \item 这是掠夺性贷款的第三个特征。贷款人通过欺骗手段隐瞒贷款的真实性质，例如将可调整利率贷款伪装成固定利率贷款，这属于掠夺性贷款行为。
    \end{itemize}
\end{explanation}

\checklist{掠夺性借贷的特征}


\begin{question}
    Which of the following was not a cause of the misalignment between investors, and rating agencies' incentives prior to the credit crisis of 2007-2009?
\end{question}

\begin{choices}
    \item Profit motive of the rating agencies
    \item Pressure from the originators of securitized products
    \item Manipulation of the ratings process by the originators
    \item Investors' lack of understanding of the products they were purchasing
\end{choices}

\begin{correctanswer}
    D
\end{correctanswer}

\begin{explanation}
    \textbf{A选项错误。}
    \begin{itemize}
    \item 评级机构的利润动机是导致激励错位的原因之一。评级机构为了获取更多业务和利润，可能会降低评级标准，这导致了评级质量的下降。
    \end{itemize}

    \textbf{B选项错误。}
    \begin{itemize}
    \item 证券化产品发起人的压力是导致激励错位的重要原因。发起人通过向评级机构施压，要求获得更高的评级，这影响了评级的客观性。
    \end{itemize}

    \textbf{C选项错误。}
    \begin{itemize}
    \item 发起人对评级过程的操纵是导致激励错位的关键因素。发起人利用对评级机构评级流程的了解，设计出能够获得AAA评级的产品，这导致了评级的失真。
    \end{itemize}

    \textbf{D选项正确。}
    \begin{itemize}
    \item 投资者对产品缺乏了解不是导致激励错位的原因。根据金融危机调查委员会的调查结果，激励错位的主要原因是评级机构的模型缺陷、利润动机、发起人压力、市场份额驱动以及缺乏有效的公共监督等。
    \end{itemize}
\end{explanation}

\checklist{评级机构激励错位的原因}

\begin{question}
    Which of the following statements is correct regarding credit risk scoring models?
\end{question}

\begin{choices}
    \item A pooled model will result in scores ranging from 300 to 850
    \item A custom model is cheaper to implement than credit bureau scores
    \item Multiple requests for new credit will reduce an applicant's credit score
    \item An example of a characteristic in a scoring model is the applicant's current gross salary of \$50,000
\end{choices}

\begin{correctanswer}
    C
\end{correctanswer}

\begin{explanation}
    \textbf{A选项错误。}
    \begin{itemize}
    \item 信用局评分模型（而不是数据池模型）的分数范围是300到850。数据池模型使用相似的外部数据或内部数据，其分数范围可能不同。
    \end{itemize}

    \textbf{B选项错误。}
    \begin{itemize}
    \item 客户模型比信用局评分模型实施成本更高。客户模型需要银行自己建立和维护，而信用局评分模型可以直接使用外部专业机构提供的服务。
    \end{itemize}

    \textbf{C选项正确。}
    \begin{itemize}
    \item 多次申请新信用会降低申请人的信用评分。信用记录会显示信用申请历史，频繁的申请会被视为负面因素，导致信用评分下降。
    \end{itemize}

    \textbf{D选项错误。}
    \begin{itemize}
    \item "当前雇主的总工资"是一个特征（characteristic），而具体的工资数额（\$50,000）是一个属性（attribute）。在评分模型中，特征是指变量类别，而属性是指具体的数值。
    \end{itemize}
\end{explanation}

\checklist{信用评分模型的特点}


\begin{question}
    At the beginning of the year, a firm bought an AA-rated corporate bond at USD 115 per USD 100 face value. Using market data, the risk manager estimates the following year-end values for the bond based on interest rate simulations informed by the economics team:

    \begin{table}[htbp]
        \centering
        \begin{tabular}{|c|c|}
            \hline
            Rating & Year-end Bond Value \\
            & (USD per USD 100 face value) \\
            \hline
            AAA     & 118 \\
            AA      & 114 \\
            A       & 108 \\
            BBB     & 103 \\
            BB      & 95  \\
            B       & 85  \\
            CCC     & 75  \\
            Default & 45  \\
            \hline
        \end{tabular}
        \caption{不同评级债券年末价值}
    \end{table}

    In addition, the risk manager estimates the 1-year transition probabilities on the AA-rated corporate bond: 
    \begin{table}[htbp]
        \centering
        \begin{tabular}{|c|c|}
            \hline
            Rating & Probability of State \\
            \hline
            AAA     & 4.00\% \\
            AA      & 82.00\% \\
            A       & 8.00\% \\
            BBB     & 4.50\% \\
            BB      & 0.40\% \\
            B       & 0.30\% \\
            CCC     & 0.20\% \\
            Default & 0.60\% \\
            \hline
        \end{tabular}
        \caption{各评级状态概率}
    \end{table}

    What is the 1-year 95\% credit VaR per USD 100 of face value closest to?
\end{question}

\begin{choices}
    \item USD 12
    \item USD 20
    \item USD 35
    \item USD 45
\end{choices}

\begin{correctanswer}
    A
\end{correctanswer}

\begin{explanation}
    \textbf{A选项正确。}
    \begin{itemize}
        \item 度量口径：信用VaR以\textbf{当前市价}为基准，取一年期末价格分布的95\%分位数对应的价格，VaR=当前价−该分位价格。
        \item 从最差到最好按价格排序并累加违约/降级概率：
        \[
        \begin{aligned}
        &\text{Default }(45):0.60\%\;\Rightarrow\;0.60\%\\
        &\text{CCC }(75):0.20\%\;\Rightarrow\;0.80\%\\
        &\text{B }(85):0.30\%\;\Rightarrow\;1.10\%\\
        &\text{BB }(95):0.40\%\;\Rightarrow\;1.50\%\\
        &\text{BBB }(103):4.50\%\;\Rightarrow\;6.00\%
        \end{aligned}
        \]
        \item 5\%左尾分位（即95\%分位价格）落在BBB状态（累计从1.5\%跳到6.0\%覆盖5\%），对应价格\$103。
        \item 当前持有价\$115，故一年期95\%信用VaR：
        \[
        \text{VaR}_{95\%}=115-103=12\;\text{(USD/100面值)}
        \]
    \end{itemize}
\end{explanation}

\checklist{信用VaR的计算方法}


\begin{question}
    A risk analyst at a bank is preparing a presentation on climate-related financial risks. The analyst focuses on the amplifiers and mitigants of climate-related financial risks and assesses the manner in which they function. Which of the following statements would be correct for the analyst to include in the presentation?
\end{question}

\begin{choices}
    \item Climate-related risks increase through asset securitizations
    \item Climate-related risks are typically impossible to hedge
    \item Physical and transition risk drivers interact with each other
    \item Transition risk drivers generally affect firms idiosyncratically
\end{choices}

\begin{correctanswer}
    C
\end{correctanswer}

\begin{explanation}
    \textbf{A选项错误。}
    \begin{itemize} 
    \item 资产证券化并不会增加气候相关风险。资产证券化是一种风险分散和转移的工具，它本身不会放大气候风险。
    \end{itemize}
    \textbf{B选项错误。}
    \begin{itemize}
    \item 气候相关风险并非完全无法对冲。虽然完全对冲可能很困难，但通过保险、衍生品等工具可以部分对冲这些风险。
    \end{itemize}

    \textbf{C选项正确。}
    \begin{itemize}
    \item 物理风险和转型风险驱动因素会相互影响。例如，气候变化导致的物理风险可能加速向低碳经济的转型，而转型政策又可能影响物理风险的程度。
    \end{itemize}

    \textbf{D选项错误。}
    \begin{itemize}
    \item 转型风险驱动因素通常会影响整个行业或经济系统，而不是仅仅影响个别公司。例如，碳定价政策会影响所有高碳排放企业。
    \end{itemize}
\end{explanation}

\checklist{气候相关风险的放大和缓解机制}





\newpage


\section{考点5:主权国家违约风险}
\setcounter{question}{0}


\begin{question}
    You are currently long \$12,000,000 par value, 7\% XYZ bonds. To hedge your position, you must decide between credit protection via a 5-year CDS with 70bp annual premiums or digital swap with 55\% payout with 60bp annual premiums. After one year, XYZ has defaulted on its debt obligations and currently trades at 65\% of par. Which of the following statements is true?
\end{question}

\begin{choices}
    \item The contingent payment from the protection buyer to the protection seller is greater under the single-name CDS than the digital swap
    \item The contingent payment from the protection buyer to the protection seller is less under the single-name CDS than the digital swap
    \item The contingent payment from the protection seller to the protection buyer is greater under the single-name CDS than the digital swap
    \item The contingent payment from the protection seller to the protection buyer is less under the single-name CDS than the digital swap
\end{choices}

\begin{correctanswer}
    D
\end{correctanswer}

\begin{explanation}
    \textbf{A选项错误。}    
    \begin{itemize}
    \item 违约时，赔付是从保护卖方支付给保护买方，而不是从买方支付给卖方。因此这个选项描述的方向是错误的。
    \end{itemize}

    \textbf{B选项错误。}
    \begin{itemize}
    \item 同样，这个选项也错误地描述了赔付方向。违约时，赔付是从保护卖方支付给保护买方。
    \end{itemize}

    \textbf{C选项错误。}
    \begin{itemize}
    \item 实物交割CDS的赔付金额是面值的35\%（1-0.65），而数字互换的赔付金额是面值的55\%。因此，数字互换的赔付金额更大。
    \end{itemize}

    \textbf{D选项正确。}
    \begin{itemize}
    \item 实物交割CDS的赔付金额是面值的35\%（1-0.65），而数字互换的赔付金额是面值的55\%。因此，实物交割CDS的赔付金额小于数字互换。
    \end{itemize}
\end{explanation}

\checklist{数字式CDS的结算方式}


\begin{question}
    A risk manager is advising the trading desk about entering into a digital credit default swap as a way to obtain credit protection. Which cash flow and delivery requirement will the desk most likely experience in the event of a default of the underlying reference asset?
\end{question}

\begin{choices}
    \item Receive the pre-agreed cash payment; delivering nothing
    \item Receive [(Par Value) – (Market Value of Reference Asset)]; deliver the reference asset
    \item Receive [(Par Value) – (Market Value of Reference Asset)]; deliver nothing
    \item Receive the pre-agreed cash payment; deliver the reference asset
\end{choices}

\begin{correctanswer}
    A
\end{correctanswer}

\begin{explanation}
    \textbf{A选项正确。}
    \begin{itemize}
    \item 数字式信用违约互换（Digital CDS）在违约时会支付预先约定的固定金额，且不需要交付任何资产。这种设计特别适用于流动性差、难以定价的参考资产。
    \end{itemize}

    \textbf{B选项错误。}
    \begin{itemize}
    \item 这个选项描述的是实物交割CDS的特点，而不是数字式CDS。实物交割CDS需要交付参考资产，并支付面值与市场价值的差额。
    \end{itemize}

    \textbf{C选项错误。}
    \begin{itemize}
    \item 这个选项描述的是现金结算CDS的特点，而不是数字式CDS。现金结算CDS支付面值与市场价值的差额，但不需要交付资产。
    \end{itemize}

    \textbf{D选项错误。}
    \begin{itemize}
    \item 数字式CDS不需要交付参考资产，只需要支付预先约定的固定金额。这个选项错误地要求交付参考资产。
    \end{itemize}
\end{explanation}

\checklist{数字式CDS的结算方式}


\begin{question}
    A default swap acts like a:
\end{question}

\begin{choices}
    \item Call option on the reference obligation for the buyer of the swap
    \item Put option on the reference obligation for the buyer of the swap
    \item Look back option on the reference obligation for the buyer of the swap
    \item Up-and-out option on the reference obligation for the buyer of the swap
\end{choices}

\begin{correctanswer}
    B
\end{correctanswer}

\begin{explanation}
    \textbf{A选项错误。}
    \begin{itemize}
    \item 看涨期权赋予买方以固定价格购买资产的权利，而违约互换是保护买方免受违约损失的工具，两者性质不同。
    \end{itemize}

    \textbf{B选项正确。}
    \begin{itemize}
    \item 违约互换的行为类似于看跌期权。当参考资产违约时，买方获得赔付，这限制了买方的下行风险，这与看跌期权的保护功能相似。
    \end{itemize}

    \textbf{C选项错误。}
    \begin{itemize}
    \item 回望期权允许买方在期权有效期内选择最优执行价格，而违约互换的赔付是预先确定的，两者机制不同。
    \end{itemize}

    \textbf{D选项错误。}
    \begin{itemize}
    \item 向上敲出期权在资产价格超过特定水平时自动失效，而违约互换在违约发生时才会触发赔付，两者触发条件不同。
    \end{itemize}
\end{explanation}

\checklist{违约互换的性质}

\begin{question}
    In comparing agency rating systems to internal (experts-based) rating systems, evidence has shown that:
\end{question}

\begin{choices}
    \item Internal systems and agency systems are equally compliant in regard to specificity
    \item Agency systems are more compliant in regard to verifiability than internal systems
    \item Agency systems are less compliant in regard to measurability than internal systems
    \item Internal systems are more compliant in regard to objectivity and homogeneity than agency systems
\end{choices}

\begin{correctanswer}
    B
\end{correctanswer}

\begin{explanation}
    \textbf{A选项错误。}
    \begin{itemize}
    \item 内部评级系统和机构评级系统在特异性方面并不等同。机构评级系统通常有更明确的评级标准和更一致的评级方法。
    \end{itemize}

    \textbf{B选项正确。}
    \begin{itemize}
    \item 机构评级系统在可验证性方面比内部评级系统更符合要求。机构评级系统有更透明的评级过程和更严格的验证机制，使得评级结果更容易被验证。
    \end{itemize}

    \textbf{C选项错误。}
    \begin{itemize}
    \item 机构评级系统在可测量性方面并不比内部评级系统差。相反，机构评级系统通常有更完善的量化指标和更系统的测量方法。
    \end{itemize}

    \textbf{D选项错误。}
    \begin{itemize}
    \item 内部评级系统在客观性和同质性方面并不比机构评级系统更好。机构评级系统通常有更严格的评级标准和更统一的评级方法，能够提供更客观和一致的评级结果。
    \end{itemize}
\end{explanation}

\checklist{机构评级系统的特点}

\begin{question}
    The buyer of a credit-default swap (CDS):
\end{question}

\begin{choices}
    \item Has no risk exposure to the combined holdings of the CDS and reference obligation
    \item Is subject to the credit risk of the seller
    \item Must place collateral with the seller of the CDS to ensure performance
    \item Can surrender the CDS for the value of the accumulated payments
\end{choices}

\begin{correctanswer}
    B
\end{correctanswer}

\begin{explanation}
    \textbf{A选项错误。}
    \begin{itemize}
    \item CDS买方仍然面临风险敞口。即使持有CDS和参考资产，买方仍然面临卖方的信用风险，无法完全消除风险。
    \end{itemize}

    \textbf{B选项正确。}
    \begin{itemize}
    \item CDS买方面临卖方的信用风险。如果参考资产违约，而卖方无法履行支付义务，买方将遭受损失。这种风险受两个因素影响：1）参考资产违约概率随时间增加，卖方违约概率也可能增加；2）经济和运营因素可能导致参考资产违约与卖方违约之间的相关性增加。
    \end{itemize}

    \textbf{C选项错误。}
    \begin{itemize}
    \item CDS买方不需要向卖方提供抵押品。相反，通常是卖方需要提供抵押品来确保其履行支付义务的能力。
    \end{itemize}

    \textbf{D选项错误。}
    \begin{itemize}
    \item CDS买方不能通过放弃CDS来获得累积付款的价值。CDS是一种信用保护工具，不能像债券那样在二级市场交易或赎回。
    \end{itemize}
\end{explanation}

\checklist{CDS买方的风险暴露}

\begin{question}
    A quantitative risk analyst at a broker-dealer bank is studying correlations between financial assets, using historical data over the last twenty years, with particular focus on the 2007-2009 global financial crisis. The analyst also examines the tools and products used to gain exposure to or hedge against changes in correlations. Which of the following can the analyst correctly conclude?
\end{question}

\begin{choices}
    \item The spreads between prices of CDO tranches will decrease as mortgage default probabilities increase during a systemic housing market crash
    \item Both VaR and ES models incorporate the impacts of dynamic financial correlations between assets to manage risk and estimate capital requirements
    \item The correlation swap is the only financial product available to market participants to hedge or expand their exposure to the correlation between assets
    \item The CDS spread would decrease if the default correlation between the reference asset and the seller of the CDS on that asset increases
\end{choices}

\begin{correctanswer}
    D
\end{correctanswer}

\begin{explanation}
    \textbf{A选项错误。}
    \begin{itemize}
    \item 在系统性房地产市场崩溃期间，随着抵押贷款违约概率增加，CDO分层的价格利差会扩大而不是缩小。这是因为违约风险增加会导致各层级的风险溢价上升。
    \end{itemize}

    \textbf{B选项错误。}
    \begin{itemize}
    \item VaR和ES模型并不总是能够充分捕捉资产之间的动态金融相关性。这些模型通常假设相关性是稳定的，而实际上相关性在危机期间会发生显著变化。
    \end{itemize}

    \textbf{C选项错误。}
    \begin{itemize}
    \item 相关性互换不是唯一可用于对冲或扩大资产相关性敞口的金融产品。市场参与者还可以使用其他工具，如CDO、CDS指数等。
    \end{itemize}

    \textbf{D选项正确。}
    \begin{itemize}
    \item 如果参考资产与CDS卖方之间的违约相关性增加，CDS利差会下降。这是因为违约相关性增加意味着卖方更可能无法履行支付义务，从而降低了CDS的价值。
    \end{itemize}
\end{explanation}

\checklist{信用相关性对CDS利差的影响}


\end{document}